\appendix

\chapter*{APÊNDICE A - Protocolo experimental consolidado}
\addcontentsline{toc}{chapter}{APÊNDICE A - Protocolo experimental consolidado}

Este apêndice consolida os elementos do protocolo experimental que sustentam os Capítulos~\ref{chap:metodo} e~\ref{chap:resultados}, como referência integrada para auditoria do trabalho.

\section*{A.1 Identificadores de execução e \textit{fingerprints}}

Cada execução é identificada de forma determinística por um \texttt{run\_id}, o SHA-256 da representação JSON canônica (chaves ordenadas) de nove campos: ativo, \textit{hash} do conjunto de \textit{features}, número do \textit{trial}, \textit{fold}, semente, versão do modelo, assinatura de configuração, assinatura das partições e versão do \textit{pipeline}. Campos opcionais ausentes participam da chave como nulos. Repetições com a mesma especificação produzem o mesmo \texttt{run\_id}; diferenças entre execuções só podem decorrer de mudança explícita em algum desses campos.

\textit{Fingerprints} complementares: a impressão do conjunto de dados resume ativo, intervalo temporal, número de linhas, somas de fechamento e de volume e o \textit{hash} do arquivo; a impressão das partições resume as listas ordenadas de sessões de \texttt{train}, \texttt{early\_stop}, \texttt{calib} e \texttt{test}; a assinatura de configuração resume hiperparâmetros, \textit{features}, tipagem, grade de quantis, horizonte máximo e id de coorte, excluindo campos voláteis como datas de execução. Comparações pareadas exigem coincidência exata da impressão das partições.

\section*{A.2 Previsões persistidas}

As previsões fora da amostra são persistidas em formato longo, com uma linha por (execução, partição, horizonte, sessão de decisão, \textit{target timestamp}, nível de quantil). Cada linha carrega o valor bruto do modelo, o valor após o rearranjo monotônico e um indicador de que o rearranjo alterou a ordem dos níveis. A grade de níveis é dado, e não estrutura: alterar a grade acrescenta linhas, sem mudar o esquema. Candidato e comparadores usam o mesmo formato e a mesma coorte.

Os dados brutos e os fatos de execução são imutáveis; toda tabela de métricas, testes e decisão é derivada deles e pode ser reconstruída sem re-treinar os modelos --- inclusive para reanálises com outra variante quantílica ou outro parâmetro de teste.

\section*{A.3 Rearranjo monotônico e \textit{gate} de degeneração}

O rearranjo ordena os \(K\) valores brutos da grade de cada previsão, produzindo quantis não decrescentes ao longo dos níveis; o indicador de rearranjo registra se a ordem mudou. Se algum valor bruto não é finito, a grade é preservada sem reordenação, e a decisão sobre sua validade fica com os critérios de qualidade. O rearranjo não altera grades colapsadas (\(\hat{q}_{\tau_1} = \hat{q}_{\tau_K}\)); estas são identificadas por um \textit{gate} de degeneração, que invalida as métricas da linha e reporta a taxa por execução. A variante (bruta ou rearranjada) que alimenta cada métrica confirmatória é fixada no pré-registro; métricas de risco usam sempre a variante rearranjada.

\section*{A.4 \textit{Scorecard} e regra de decisão}

A versão completa do \textit{scorecard} está na Seção~\ref{sec:scorecard}. Em síntese: (1) \textit{gate} de validade estatística --- perda \textit{pinball}, sem seleção dos melhores \textit{runs}, família de Holm restrita às mesmas partições, série deduplicada e alinhada; (2) \textit{gate} de calibração (H1) por horizonte, contra bandas pré-registradas; (3) veredito de H2 por horizonte, via DM/Holm contra cada nível da hierarquia de \textit{baselines} e pertencimento ao MCS; (4) H3 por consistência em pelo menos dois de três métodos. O veredito só é emitido com todas as condições de validade satisfeitas, e as demais métricas nunca o alteram.

\section*{A.5 Protocolo estatístico}

A comparação fora da amostra exige interseção exata por (ativo, partição, horizonte, \textit{target timestamp}, nível de quantil), com uma única observação por ponto alinhado (a previsão operacionalmente mais recente). O teste primário é o de Diebold-Mariano \cite{DIEBOLDMARIANO1995}, unilateral, sobre a \textit{pinball loss}, com variância HAC de defasagem \(h-1\) \cite{NEWEYWEST1987} e correção HLN \cite{HARVEY1997}; o ajuste de múltiplas comparações é o de Holm \cite{HOLM1979}; o \textit{Model Confidence Set} \cite{HANSEN2011MCS} usa \textit{bootstrap} em blocos com parâmetros fixados no pré-registro. Todos os procedimentos são conferidos contra implementações de referência. Toda comparação confirmatória é restrita a execuções da mesma coorte.

\section*{A.6 Pré-registro}

A rodada confirmatória só se inicia após o congelamento do pré-registro, que fixa: hipóteses (H1, H2, H3); candidato (\textit{features}, hiperparâmetros, sementes, \textit{folds}); grade de quantis e horizontes; hierarquia de \textit{baselines} e suas convenções de emissão; métrica primária; bandas de calibração e pares de níveis dos intervalos; variante quantílica de cada métrica; parâmetros do DM, do Holm e do MCS; variante da predição conformal; limiar de degeneração; e a regra de decisão do \textit{scorecard}. O documento é \textit{hasheado}, e o \textit{hash} é ancorado aos artefatos confirmatórios: qualquer alteração posterior é detectável.

\chapter*{APÊNDICE B - Glossário técnico de termos recorrentes}
\addcontentsline{toc}{chapter}{APÊNDICE B - Glossário técnico de termos recorrentes}

Este apêndice reúne, em formato sintético, os termos técnicos recorrentes ao longo do trabalho.

\section*{B.1 Termos de previsão probabilística}

\phantomsection\label{gl:calibracao}\textbf{Calibração probabilística}: correspondência entre a probabilidade nominal declarada pelo modelo e a frequência empírica observada. Um modelo bem calibrado no nível \(0{,}9\) tem o valor observado abaixo do quantil previsto em aproximadamente \(90\%\) dos casos.

\phantomsection\label{gl:nitidez}\textbf{Nitidez (\textit{sharpness})}: concentração da distribuição preditiva; intervalos mais estreitos são mais nítidos. Só é desejável sujeita à calibração.

\phantomsection\label{gl:grade}\textbf{Grade de quantis}: conjunto comum de níveis \(\{\tau_1 < \dots < \tau_K\}\), com \(K\) entre 7 e 9, emitido por todos os modelos.

\phantomsection\label{gl:pinball}\textbf{\textit{Pinball loss}}: perda assimétrica da regressão quantílica, definida na Equação~\ref{eq:pinball}; estritamente consistente para o quantil do nível correspondente e métrica primária deste trabalho.

\phantomsection\label{gl:crps}\textbf{CRPS (\textit{Continuous Ranked Probability Score})}: regra de pontuação própria para a distribuição preditiva inteira; igual a duas vezes a integral da \textit{pinball} sobre os níveis (Equação~\ref{eq:crps}). Métrica complementar.

\phantomsection\label{gl:picp}\textbf{PICP (\textit{Prediction Interval Coverage Probability})}: frequência empírica de observações dentro do intervalo delimitado por um par de níveis \((\tau_l, \tau_h)\) (Equação~\ref{eq:picp}), comparada ao nominal \(\tau_h - \tau_l\).

\phantomsection\label{gl:mpiw}\textbf{MPIW (\textit{Mean Prediction Interval Width})}: largura média do intervalo de predição (Equação~\ref{eq:mpiw}); sempre lida em conjunto com o PICP.

\phantomsection\label{gl:crossing}\textbf{\textit{Quantile crossing}}: violação da ordenação dos quantis previstos (\(\hat{q}_{\tau_j} > \hat{q}_{\tau_k}\) com \(\tau_j < \tau_k\)), típica de regressores treinados por nível.

\phantomsection\label{gl:rearranjo}\textbf{Rearranjo monotônico}: ordenação dos valores da grade de cada previsão, que restaura a monotonicidade sem afastar a previsão da curva quantílica verdadeira \cite{CHERNOZHUKOV2010REARRANGEMENT}; corrige ordem, não calibração.

\phantomsection\label{gl:degeneracao}\textbf{Degeneração da grade}: previsão cujos níveis extremos coincidem, sem informação distribucional; tratada por \textit{gate} próprio, separado do rearranjo.

\phantomsection\label{gl:cqr}\textbf{CQR (\textit{Conformalized Quantile Regression})}: ajuste conformal dos quantis de um modelo com base nos escores de não conformidade de um conjunto de calibração disjunto; usada como referência de cobertura empírica.

\section*{B.2 Termos de protocolo e governança}

\phantomsection\label{gl:leakage}\textbf{\textit{Leakage} temporal}: uso de informação futura no cálculo de \textit{features}, no ajuste de transformações ou na seleção de modelo. Invalida conclusões fora da amostra.

\phantomsection\label{gl:walkforward}\textbf{\textit{Walk-forward} com purga e embargo}: validação temporal com origens sucessivas, em que o treino cresce a cada \textit{fold} e as partições são separadas por \(h_{\max}\) + embargo sessões excluídas.

\phantomsection\label{gl:split}\textbf{Partições do \textit{fold}}: \texttt{train} (ajuste), \texttt{early\_stop} (parada antecipada e avaliação da busca exploratória), \texttt{calib} (calibração conformal, nunca usada para seleção) e \texttt{test} (avaliação fora da amostra), disjuntas e nessa ordem temporal.

\phantomsection\label{gl:known}\textbf{Entradas conhecidas e observadas}: no TFT, conhecidas são as predeterminadas até o horizonte (calendário); observadas são as medidas apenas até a decisão (preço, indicadores, sentimento, fundamentos).

\phantomsection\label{gl:coorte}\textbf{Coorte}: conjunto de execuções com mesmo ativo, conjunto de \textit{features}, horizonte máximo, geometria de partições e id de coorte. Comparações confirmatórias ocorrem apenas dentro de uma coorte.

\phantomsection\label{gl:dedup}\textbf{Deduplicação \textit{operationally-latest}}: quando \textit{folds} distintos preveem o mesmo ponto alinhado, mantém-se apenas a previsão mais recente, garantindo uma observação por ponto.

\phantomsection\label{gl:preregistro}\textbf{Pré-registro}: documento congelado e \textit{hasheado} antes da rodada confirmatória, que fixa hipóteses, candidato, \textit{baselines}, métricas, bandas, parâmetros dos testes e regra de decisão.

\phantomsection\label{gl:scorecard}\textbf{\textit{Scorecard}}: aplicação mecânica da regra de decisão pré-registrada, que separa o veredito do perfil comparativo.

\section*{B.3 Termos do modelo}

\phantomsection\label{gl:tft}\textbf{TFT (\textit{Temporal Fusion Transformer})}: modelo de previsão multi-horizonte que combina seleção de variáveis, camadas recorrentes, atenção \textit{multi-head} e saída quantílica \cite{LIMETAL2021TFT}.

\phantomsection\label{gl:vsn}\textbf{VSN (\textit{Variable Selection Network})}: subcomponente do TFT que atribui pesos por passo de tempo às variáveis de entrada.

\phantomsection\label{gl:feature-families}\textbf{Famílias de \textit{features}}: preço (OHLCV e derivadas de preço e volume), técnico (indicadores e medidas de volatilidade), sentimento (agregados de notícias por sessão) e fundamento (demonstrativos e razões contábeis).

\phantomsection\label{gl:gbm}\textbf{GBM quantílico}: \textit{gradient boosting} treinado com a \textit{pinball loss}, um modelo por par (nível, horizonte); comparador forte da hierarquia de \textit{baselines}.

\phantomsection\label{gl:run-id}\textbf{\texttt{run\_id}}: identidade determinística de uma execução, conforme o Apêndice A.1.

\section*{B.4 Termos de teste estatístico}

\phantomsection\label{gl:dm}\textbf{DM (Diebold-Mariano)}: teste de igualdade de acurácia preditiva esperada entre dois modelos, sobre a série de perdas pareadas por \textit{target timestamp}.

\phantomsection\label{gl:hln}\textbf{HLN (Harvey-Leybourne-Newbold)}: correção do DM para amostras pequenas, com distribuição \(t\) de Student.

\phantomsection\label{gl:holm}\textbf{Holm}: ajuste sequencial de múltiplas comparações que controla a \textit{Family-Wise Error Rate}.

\phantomsection\label{gl:mcs}\textbf{MCS (\textit{Model Confidence Set})}: procedimento de \cite{HANSEN2011MCS} que identifica o conjunto de modelos não rejeitados como superiores ao nível \(\alpha\).

\phantomsection\label{gl:christoffersen}\textbf{Christoffersen e Kupiec}: testes de razão de verossimilhança para a cobertura de intervalos e VaR; Kupiec avalia a proporção de violações e Christoffersen acrescenta a independência entre elas.

\chapter*{APÊNDICE C - Reprodução}
\addcontentsline{toc}{chapter}{APÊNDICE C - Reprodução}

Este apêndice indica como reproduzir o \textit{pipeline}. O repositório do trabalho contém o código, a configuração e a documentação de decisões.

\section*{C.1 Ambiente}

O ambiente é definido por uma imagem Docker com Python 3.12 e dependências fixadas em arquivo de travamento. Os alvos principais são \texttt{make setup} (configuração inicial), \texttt{make check} (verificações estáticas e testes) e \texttt{make test-cov} (testes com relatório de cobertura); \texttt{make docker-build}, \texttt{make docker-up} e \texttt{make docker-shell} constroem e acessam o ambiente em contêiner.

\section*{C.2 Dados e \textit{features}}

As fontes são \texttt{yfinance} (cotações diárias), Alpha Vantage (notícias, demonstrativos e datas de divulgação) e o modelo FinBERT com revisão de pesos fixada. A sequência é: ingestão de cotações, notícias e fundamentos; pontuação e agregação de sentimento por sessão; e construção do conjunto de dados com as 55 \textit{features}, o alvo e o \textit{gate} de qualidade.

\section*{C.3 Treinamento e avaliação}

Sobre o conjunto de dados, executam-se: os cinco \textit{baselines} estatísticos, o GBM quantílico e o TFT sob o mesmo protocolo \textit{walk-forward}; a busca exploratória de hiperparâmetros do TFT; e a rodada confirmatória do cohort pré-registrado. Em seguida, calculam-se métricas, testes, predição conformal e contribuição de famílias, e o \textit{scorecard} produz o veredito. As execuções são registradas em MLflow com \textit{backend} local.

\chapter*{APÊNDICE D - Notas de implementação}
\addcontentsline{toc}{chapter}{APÊNDICE D - Notas de implementação}

Este apêndice descreve, de forma sucinta, as escolhas de engenharia de software que sustentam o trabalho. Elas não fazem parte do objeto científico, mas asseguram que os números reportados sejam verificáveis.

\section*{D.1 Arquitetura}

O código segue arquitetura hexagonal (\textit{ports and adapters}) organizada em \textit{vertical slices} por contexto: dados de mercado, engenharia de \textit{features}, armazenamento analítico e modelagem, com os contextos de avaliação e inferência previstos na mesma estrutura. O domínio é implementado apenas com a biblioteca padrão do Python; bibliotecas numéricas e de aprendizado de máquina (\texttt{pandas}, \texttt{pyarrow}, \texttt{torch}, \texttt{pytorch-forecasting}, \texttt{lightgbm}, \texttt{statsforecast}, \texttt{optuna}) ficam restritas a adaptadores. A direção das dependências é verificada automaticamente pelo \texttt{import-linter}, com checagem de tipos estrita (\texttt{mypy --strict}) e cobertura mínima de testes de 90\% como condição para integrar qualquer mudança.

\section*{D.2 Validação das estatísticas}

As métricas e os testes são implementados como funções puras sobre objetos de valor com invariantes (grade de quantis, séries de perdas pareadas, séries de cobertura) e validados por testes de contrato contra casos analíticos e contra implementações de referência (R, \texttt{scikit-learn}, \texttt{scoringrules}, \texttt{arch}, \texttt{statsmodels}). Cada indicador técnico é conferido contra sua fórmula canônica e contra um teste de ausência de \textit{leakage}.

\section*{D.3 Armazenamento}

Os dados seguem a organização em camadas \textit{bronze} (dados brutos imutáveis), \textit{silver} (fatos de execução, fonte da verdade) e \textit{gold} (tabelas analíticas derivadas e reconstruíveis), em arquivos Parquet consultados com DuckDB. A camada \textit{silver} contém as tabelas de dimensão de execuções, configuração, previsões fora da amostra (formato longo, com nível de quantil na chave), métricas por partição e falhas; fatos são apenas acrescentados, nunca sobrescritos. O rastreamento de experimentos usa MLflow com \textit{backend} SQLite local.

\section*{D.4 Registro de decisões}

As decisões de desenho são registradas em documentos de decisão arquitetural (ADR) versionados no repositório, por exemplo: enquadramento em calibração probabilística (ADR 0.0.0002), controle de \textit{leakage} (0.0.0018), convenções de emissão de quantis dos \textit{baselines} (0.0.0052), definição do alvo (3.5.0001), formato longo das previsões (4.1.0002), indexação do \textit{target timestamp} por sessão (4.3.0001), rearranjo sobre a grade densa (4.3.0002), \textit{walk-forward} expansivo e partição de calibração (5.1.0001--5.1.0003), GBM direto por nível e horizonte (5.3.0001) e configuração do TFT (5.4.0001--5.4.0006).
